% VI36-DETAILED-PROBLEM-18
\begin{problem}[VI.36.P18: Tensoring the basic twists]\label{prob:vi36-18}
Prove directly from transition factors that
\[
\mathcal O(d)\otimes\mathcal O(e)\cong\mathcal O(d+e)
\]
on \(\mathbb P^n_k\).

\textbf{Provenance.} Canonical VI/36 extension.
\end{problem}

\ifdefined\FullProblemDossiers
\begin{solution}
On \(U_i\), let
\[
e_i^{(d)}=x_i^d,
\qquad
e_i^{(e)}=x_i^e.
\]
Their tensor product gives a local frame
\[
e_i^{(d)}\otimes e_i^{(e)}.
\]
Map it to
\[
e_i^{(d+e)}=x_i^{d+e}.
\]
This defines an isomorphism of free rank-one modules on each chart.

On an overlap,
\[
e_i^{(d)}
=
\left(\frac{x_i}{x_j}\right)^d e_j^{(d)},
\]
and
\[
e_i^{(e)}
=
\left(\frac{x_i}{x_j}\right)^e e_j^{(e)}.
\]
Therefore the tensor frame transforms by
\[
\left(\frac{x_i}{x_j}\right)^{d+e},
\]
which is exactly the transition factor of \(\mathcal O(d+e)\). Hence the local identifications agree on overlaps and
glue:
\[
\boxed{
\mathcal O(d)\otimes\mathcal O(e)\cong\mathcal O(d+e).
}
\]
This is a concrete identity for the basic twists, not yet the general Picard-group theory.
\end{solution}
\fi
